\documentclass[10pt]{article}

\usepackage[letterpaper, portrait, margin=1.25in]{geometry}

\usepackage{authblk}
\usepackage[yyyymmdd,hhmmss]{datetime}
\usepackage{fancyhdr}
\pagestyle{fancy}
\lhead{ASTR 324: Introduction to AstroStatistics and Big Data in Astronomy}
\rhead{Spring 2018}
\rfoot{\em \tiny Compiled on \today\ at \currenttime}
\cfoot{\thepage}
\lfoot{}


%%%%%%%%%%%%%%%%%%%%%%%%%%%%%%%%%%%%%%%%%%%%%%%%%%%%
%%% author-defined commands
\newcommand\about     {\hbox{$\sim$}}
\newcommand\x         {\hbox{$\times$}}
\newcommand\othername {\hbox{$\dots$}}
\def\eq#1{\begin{equation} #1 \end{equation}}
\def\eqarray#1{\begin{eqnarray} #1 \end{eqnarray}}
\def\eqarraylet#1{\begin{mathletters}\begin{eqnarray} #1 %
                  \end{eqnarray}\end{mathletters}}
\def\non    {\nonumber \\}
\def\DS     {\displaystyle}
\def\E#1{\hbox{$10^{#1}$}}
\def\sub#1{_{\rm #1}}
\def\case#1/#2{\hbox{$\frac{#1}{#2}$}}
\def\about  {\hbox{$\sim$}}
\def\x      {\hbox{$\times$}}
\def\ug               {\hbox{$u-g$}}
\def\gr               {\hbox{$g-r$}}
\def\ri               {\hbox{$r-i$}}
\def\iz               {\hbox{$i-z$}}
\def\a                {\hbox{$a^*$}}
\def\O                {\hbox{$O$}}
\def\E                {\hbox{$E$}}
\def\Oa               {\hbox{$O_a$}}
\def\Ea               {\hbox{$E_a$}}
\def\Jg               {\hbox{$J_g$}}
\def\Fg               {\hbox{$F_g$}}
\def\J                {\hbox{$J$}}
\def\F                {\hbox{$F$}}
\def\N                {\hbox{$N$}}
\def\dd               {\hbox{deg/day}}
\def\mic              {\hbox{$\mu{\rm m}$}}
\def\Mo{\hbox{$M_{\odot}$}}
\def\Lo{\hbox{$L_{\odot}$}}
\def\comm#1           {\tt #1}
\def\refto#1          {\ref #1}
\def\T#1              {({\bf #1})}
\def\H#1              {({\it #1})}

%%%%%%%%%%%%%%%%%%%%%%%%%%%%%%%%%%%%%%%%%%%%%%%%%%%%


\title{ASTR 324: AstroStatistics}
%\author{\v{Z}eljko Ivezi\'{c}  \& Mario Juri\'{c}}
\author{\v{Z}eljko Ivezi\'{c}}
\affil{University of Washington, Spring Quarter 2026}
\date{\vspace{-5ex}}

\begin{document}
\maketitle

\vskip 0.3in
\leftline{{\bf  Location and Time:} Tuesdays and Thursdays 10:00-11:20am, {\bf PAA 210}}

\vskip 0.2in
\leftline{{\bf  Office Hours:} { Any time when my office doors are open; }}
\leftline{\hskip 1.15in After Thu class, or Tue afternoon are the best.}
\vskip 0.2in
\leftline{{\bf  Grading:} 10 homeworks, 8\% each; final exam: 20\%}
\leftline{                 \hskip 0.8in key: $>$90\%=A, $>$80\%=B, $>$70\%=C, $>$50\%=D.}
\vskip 0.2in
\leftline{{\bf  Class web site:}   https://github.com/uw-astr-324/astr-324-s26}
\vskip 0.2in
\leftline{{\bf  Reference textbook:} (not required)}
\leftline{Ivezi\'{c}, Connolly, VanderPlas \& Gray: {\it Statistics, Data Mining, and Machine Learning in Astronomy: }}
\leftline{{\it A Practical Python Guide for the Analysis of Survey Data} (Princeton University Press, 2014)}
\leftline{See http://press.princeton.edu/titles/10159.html}
\vskip 0.3in

{\bf Learning Goals:}

This course will introduce students to most common statistical and computer science methods 
used in astronomy and other physical sciences. It will combine theoretical background with 
examples of data analysis based on modern astronomical datasets. Practical data analysis 
will be done using python tools, with emphasis on astroML module (see www.astroML.org). 
While focused on astronomy, this course should be useful to all students interested in data 
analysis in physical sciences and engineering. The lectures will be aimed at undergraduate 
students and the main discussion topics will be based on  Chapters 4 and 5, and selected 
topics from Chapters 6-10, from the reference textbook. 

By taking this course, students will develop basic understanding of topics such as robust 
statistics, hypothesis testing, maximum likelihood analysis, Bayesian statistics, model 
parameter estimation, the goodness of fit and model selection, density estimation and 
clustering, unsupervised and supervised classification, dimensionality reduction, 
regression and time series analysis. Most of these topics will be applied in class homeworks 
to analysis of astronomical data. 

{\bf Prerequisites:}

The students taking this class are required to have basic calculus and basic python skills, 
as well as basic scientific measurements and statistics skills at the level of a freshman lab. 

{\bf Lecture format:}

New material will typically be covered during the first class in a week, while
the second class in a week will be more focused on practical data analysis work. 

\newpage 
% \vskip 0.2in
\leftline{\bf  Class Schedule:}
\begin{itemize}
\item WEEK 1 (Mar 31/Apr 2): {\bf Overview of the course and
    introduction to statistics.} 
      Introduction to class: syllabus, literature, astroML, python,
      jupyter notebooks, GitHub. 
   Basics about statistics.
   Distributions and Random samples.
   Robust statistics.

\item WEEK 2 (Apr 7/Apr 9): {\bf Maximum Likelihood Estimation method}
What likelihood is and what it is good for?
Simple examples of MLE: one-dimensional Gaussian.
MLE in action: fitting a parametrized model with heteroscedastic gaussian errors on y axis.
Goodness of fit.
Cost functions and penalized likelihood.
Non-gaussian likelihood: binomial distribution (coin flip problem).
Conceptual difficulties with the MLE (example: "waiting for a bus" problem).
 
\item WEEK 3 (Apr 14/Apr 16): {\bf Introduction to Bayesian statistics and inference}
Bayes Rule extended to Bayesian Inference.
The role of priors in Bayesian Inference.
A simple parameter estimation example.
Nuisance parameters and marginalization.


\item WEEK 4 (Apr 21/Apr 23): {\bf Applications of Bayesian statistics and inference}
  % Apr 23: OSU colloquium 
Simple parameter estimation examples.
Bayesian model selection.
Simple model selection examples.
For overachievers: ABC and Hierarchical Bayes. {\it No in-class lecture on Apr 23.}

\item WEEK 5 (Apr 28/Apr 30): {\bf Introduction to Markov Chain Monte Carlo}
      % Apr 30 remote (Paris) 
   Introduction to MCMC.
   Non-linear Regression with MCMC. {\it No in-class lecture on Apr 30.}

\item WEEK 6 (May 5/May 7): {\bf Bayesian model selection with MCMC}
  Model selection example: finding bursts in time series.
Bayesian Blocks Algorithm.


\item WEEK 7 (May 12/May 14): {\bf Density Estimation}
  Searching for Structure in 1-D Point Data.
Gaussian Mixture Models.
Extreme Deconvolution.

\item WEEK 8 (May 19/May 21): {\bf Clustering (Unsupervised Classification)}
 % May 21: Caltech colloquium 
Introduction to Clustering. 
K-means clustering algorithm.
Clustering with Gaussian Mixture models (GMM).
Hierarchical clustering algorithm. {\it No in-class lecture on May 21.}


\item WEEK 9 (May 26/May 28): {\bf Supervised Classification}
Introduction to Supervised Classification.
An example of a discriminative classifier: Support Vector Machine classifier.
An example of a generative classifier: star/galaxy separation using Gaussian Naive Bayes classifier.
Comparison of many methods using ROC curves.

 \item WEEK 10 (June 2/June 4): {\bf Dimensionality Reduction}
The curse of dimensionality.
Principal Component Analysis.
Comparing PCA, NMF and ICA. 

\item {\bf FINAL EXAM: June 8 (Monday, 10:30am-12:20pm, PAA 210):} closed book final exam.

\end{itemize}


\vskip 0.2in

\leftline{\bf  Homework}

There will be ten homeworks, assigned on Thursdays, and due next Thursday. They will
be centered on practical work using python and designed to test the weekly progress. 

We will use modern software engineering tools, GitHub and Jupyter notebooks, for HW submission. 

\end{document}






% Week-1-Tue.key 
%     intro to class, big data, LSST and astroML 
% Week-1-Thu.ipynb
%     setup: github, class repo, astroML repo, 
\item WEEK  2 (starting Apr 3): Introduction to statistics (probability, distributions, 
             robust statistics, Central Limit Theorem,  hypothesis testing).
% Week-2-Tue.ipynb
%     stats intro, estimators, moments 
% Week-2-Thu.ipynb
%     CLT, multi-variate distributions, chi2, intro to histograms 
\item WEEK  3 (starting Apr 10):  Maximum likelihood and applications in astronomy.
% Week-3-Tue.ipynb
%     MLE, homoscedastic Gaussian, goodness-of-fit  
% Week-3-Thu.ipynb
%     LSQ, fitting straight line, optimal photometry 
\item WEEK  4 (starting Apr 17):  Bayesian statistics and introduction to MCMC.
% Week-4-Tue.ipynb
%     introduction to Bayesian inference, priors, heteroscedastic Gaussian, 
% Week-4-Thu.ipynb
%     Bayesian model selection, binomial distribution, Bayes factor, BIC, Knuth's hist., Bayesian Blocks 
\item WEEK  5 (starting Apr 24):  Model parameter estimation and model selection.
% Week-5-Tue.ipynb
%     Parameter estimation with MCMC
% Week-5-Thu.ipynb
%     Model selection using MCMC

% --------------------------------------------------------------------
\item WEEK  6 (starting May 1):   Time series analysis.
% Week-6-Tue.ipynb
%     into, Fourier analysis, gatspy, Lomb-Scargle periodogram
% Week-6-Thu.ipynb
%     stochastic time series, DRW, power-law PSD
%    HW5 <HW6 from s17-private (LINEAR periodic variables)

\item WEEK  7 (starting May 8): Dimensionality reduction and regression.
% Week-7-Tue.key 
%          Dimensionality reduction and regression 
% Week-7-Thu.key 
% Week-7-Thu.ipynb:  total LSQ, linear basis function regression 
%      HW6  <- HW8 from s17-private  (PCA in 4D LINEAR-based space, age-color regression) 

\item WEEK  8 (starting May 15):  review and Gwen 



\item WEEK  9 (starting May 22):  Density estimation and clustering.

% Week-9-Tue.key  -> Density Estimation. I    (24 pages) 
%      review of histograms, Bayesian Blocks, GMMs, 1-D KDE, data smoothing 
%   Week-9-Tue.ipynb
%        1-D GMM examples 
%        multi-D GMM examples with LINEAR  

% Week-9-Thu.key  -> Density Estimation. II   (17 pages) 
%      2D KDE, Bayesian NN, parametric density estimation, GMM and EM method, XD, 39 pages
%    Week-9-Thu.ipynb 
%              GMM clustering for sin(i) vs. a for asteroids  (same as Week-10-Tue.ipynb !!!) 
%                        better to find other examples of multi-D clustering, e.g. XD !!!      
%         %%%  perhaps 1D simulation with 3 gaussians convolved with errors and then XD? 


% in grad version from W17, Andy's notebooks:
% Week-8-Tue.ipynb   Density estimation
% Week-8-Thu.ipynb   Classification 

% HW 7 is ready!!!   KDE and GMM for LINEAR data (and  no other homeworks needed!) 


\item WEEK  10 (starting May 29):  Unsupervised and supervised classification.
% Week-10-Tue.key   (17 pages)    A BIT OF REPEAT OF WEEK 9 - a review?  
%      	•	Clustering (unsupervised classification)  
%                   – 1-D hypothesis testing                                                      
%                   – clustering with Gaussian Mixture models (GMM)    
%                   – K-means clustering algorithm                     
%                   – hierarchical clustering algorithm 
% Week-10-Tue.ipynb      
%              GMM clustering for sin(i) vs. a for asteroids  and hierarchical clustering for asteroids 
%                      (very short, find another example, perhaps LINEAR? see Thu notebook) 

% Week-10-Thu.key   (18 pages) 
%	•	(supervised) Classification        
%                   – potpourri of supervised classification methods: naive Bayes, 
%                           quadratic discriminant analysis, GMM, KNN, support vector machines
%                   – classification comparison with ROC curves
% Week-10-Thu.ipynb     (expand?) 
%            supervised classification for LINEAR variables and ROC curves 

 * Lecture 17: [Unsupervised Classification](Week-10-Tue.key) [(notebook)](Week-10-Tue.ipynb)
 * Lecture 18: [Supervised Classification](Week-10-Thu.key) [(notebook)](Week-10-Thu.ipynb)



\item FINAL EXAM: June 6 (Tue, 2:20pm, B360): cake and closed book final exam.

\end{itemize}


\vskip 0.2in

\leftline{\bf  Homework}

There will be ten homeworks, assigned on Thursdays, and due next Thursday. They will
be centered on practical work using python and designed to test the weekly progress. 

We will use modern software engineering tools, GitHub and Jupyter notebooks for HW submission. 

\end{document}

